\documentclass[10pt,a4paper]{amsart}
\usepackage[margin=19mm]{geometry}
\usepackage{amsmath,amssymb,mathtools}
\usepackage[scr=rsfs,cal=euler]{mathalfa}
\usepackage{xfrac}
\usepackage[unicode,hidelinks,bookmarks=false]{hyperref}
\newcommand{\ie}{i.e.\ }
\newcommand{\cf}{cf.\ }
\newcommand{\resp}{respectively\ }
\newcommand{\Subsection}{\subsection}
\providecommand{\nobreakdash}{\nobreak\hbox{-}}
\setlength{\emergencystretch}{2em}
\multlinegap0pt
\usepackage{amssymb}
\newtheorem{proposition}{Proposition}
\newcommand{\BLUE}[1]{{\color{blue} #1}}
\newcommand{\OLDPROOF}[1]{}
\renewcommand{\Pr}{\mathbf{P}}
\newcommand{\RED}[1]{{\color{red} #1}}
\title{``Game, Set, Match'' \\ Double Delight Watching a Grand Slam Tennis Match}
\author{Edsel A. Pena, Dip Das, Yuexuan Wu}
\date{Source: arXiv:2603.02360v1 (CC BY 4.0)}
\begin{document}
\maketitle

\noindent\emph{Source and licence.} ``Game, Set, Match'' \\ Double Delight Watching a Grand Slam Tennis Match. Version arXiv:2603.02360v1 (CC BY 4.0), distributed by arXiv under the Creative Commons Attribution 4.0 International licence, http://creativecommons.org/licenses/by/4.0/, as recorded in the arXiv licence metadata for this exact version and shown on the abstract page https://arxiv.org/abs/2603.02360v1. Full licence notice: \texttt{licenses/arxiv-2603.02360-CC-BY-4.0.txt}.\par\smallskip

\noindent\textit{Editorial scope.} Counted unit: the article's proposition probs (source0.tex lines 615--655). The article's complete original proof of it is delivered with it (source0.tex lines 657--809, one \texttt{proof} environment of 153 lines). No mathematics is written, repaired or re-derived: the statement and the proof are the article's own lines, in the article's own notation, and no retained line is substituted or re-flowed.  No further source-local context is needed: every cross-reference of every retained line resolves inside the two delivered spans.  Nothing was omitted from any retained span, so the package carries no omission record.  No retained line contains a citation, so the wrapper carries no bibliography and no bibliography entry.  No retained line contains a figure environment, an image-inclusion command or drawing code, so no image is delivered and the package declares no asset dependency.  Editorial formatting only, in addition: an AMS \texttt{amsart} preamble that restates the article's own declarations and macros used by the retained lines, namely the environments \texttt{proposition} on separate counters, together with the standard packages those lines need.  The retained environments print in this document as Proposition 1 in order of appearance, whereas the article prints the same items with the numbering of its own full document.  The complete native source of the article as distributed in the arXiv e-print for this exact version is bundled unchanged as \texttt{sources/195614/source0.tex}.
\begin{proposition}
\label{prop: set tie-breaker probs}
    If player A serves first in the set tie-breaker, the probability that he/she wins the set tie-breaker is
    %
    \begin{eqnarray*}
    \theta_{ST}(p_A,p_B,K)  & = &  \sum_{h=0}^{K-2} \theta_{ST}(K,h;p_A,p_B,K) + 
     \theta_{ST}(K-1,K-1;p_A,p_B,K) \times \\ &&
     [\theta_{STT}(p_A,p_B)]^{I_A(2K-1)} [1-\theta_{STT}(p_B,p_A)]^{1-I_A(2K-1)},
    \end{eqnarray*}
    %
    where
    %
    \begin{eqnarray*}
   \lefteqn{     \theta_{ST}(K,h;p_A,p_B,K) \equiv  \Pr\{\mbox{$A$ wins on a $(K,h)$ score}\}  } \\ & = & \left\{
        \sum_{j=0}^{K-1} {s_A(K+h-1) \choose j} {s_B(K+h-1) \choose {(K-1) - j}} \times \right. \\
        && \left.
        \left[p_A^j q_A^{s_A(K+h-1) - j}\right]
        \left[q_B^{(K-1)-j} p_B^{s_B(K+h-1) - ((K-1)-j)}\right]
        \right\}  \left[p_A^{I_A(K+h)} q_B^{1-I_A(K+h)}\right];
    \end{eqnarray*}
    %
    \begin{eqnarray*}
   \lefteqn{  \theta_{ST}(K-1,K-1;p_A,p_B,K)  \equiv   \Pr\{(K-1,K-1)\ \mbox{score}\} } \\ & = &
       \left\{
        \sum_{j=0}^{K-1} {{K-1} \choose j} {{K-1} \choose {j}}
        \left[p_A^j q_A^{(K-1) - j}\right]
        \left[q_B^{(K-1)-j} p_B^{j}\right]
        \right\}.
    \end{eqnarray*}
    %
    Alternatively, with $B(K-1,p_A)$ and $B(K-1,q_B)$ independent binomial random variables,
    %
    \begin{eqnarray*}
        \sum_{h=0}^{K-2} \theta_{ST}(K,h;p_A,p_B,K) =
        \Pr\{B(K-1,p_A) + B(K-1,q_B) \ge K\}.
    \end{eqnarray*}
    %
In addition,
    $\theta_{ST}(p_A,p_B,K) = \theta_{ST}(q_B,q_A,K).$
    %
\end{proposition}

\begin{proof}


%\BLUE{(new proof; Christmas Break)

For $h \in \{0,1,\ldots,K-2\}$, define the event
%
$$F_h = \{\mbox{$A$ wins with a score of $(K,h)$}\} \quad \mbox{and} \quad 
T_{2(K-1)} = \{\mbox{score tied at $(K-1:K-1)$\}.}$$
%
Then
%
\begin{eqnarray*}
\theta_{ST}(p_A,p_B,K) & = & \sum_{h=0}^{K-2} \Pr\{F_h|(p_A,p_B,K)\} + \\ && 
\Pr\{T_{2(K-1)}|(p_A,p_B,K)\} \Pr\{\mbox{$A$ wins STT}|(p_A,p_B,K)\}.
\end{eqnarray*}
%
For $A$ to win with a score of $(K,h)$, he/she must win $(K-1)$ points in the first $(K+h-1)$ points played, for which he/she will be serving a total of $s_A(K+h-1)$ of these points, and then he/she must also win the $(K+h)$th point played. That is,
%
\begin{eqnarray*}
F_h & = & \left\{\sum_{i=1}^{s_A(K+h-1)} V_i + \sum_{j=1}^{s_B(K+h-1)} (1 - W_j) = K-1\right\} \cap \\ &&
\left\{(I_A(K+h) = 1, V_{s_A(K+h-1)+1} = 1) \cup 
(I_A(K+h) = 0, 1-W_{s_B(K+h-1)+1} = 1)\right\}.
\end{eqnarray*}
%
The probability of this event is
%
\begin{eqnarray*} 
\lefteqn{\Pr\{F_h|(p_A,p_B,K)\}  \equiv \theta_{ST}(K,h;p_A,p_B,K) } \\
& = & \Pr\{B(s_A(K+h-1),p_A) + B(s_B(K+h-1),q_B) = K-1\} p_A^{I_A(K+h)} q_B^{1-I_A(K+h)},
\end{eqnarray*}
%
where the binomial random variables $B(s_A(K+h-1),p_A)$ and $B(s_B(K+h-1),q_B)$ are independent. Thus,
%
\begin{eqnarray*}
\lefteqn{ \theta_{ST}(K,h;p_A,p_B,K)  =   
\left[\sum_{l=0}^{K-1}
{{s_A(K+h-1)} \choose l} p_A^l q_A^{{s_A(K+h-1)}-l} \times \right. } \\ && \left. {{s_B(K+h-1)} \choose {K-1-l}} q_B^{K-1-l} p_B^{s_B(K+h-1)-(K-1-l)}\right]  p_A^{I_A(K+h)} q_B^{1-I_A(K+h)}.
\end{eqnarray*}
%
Similarly, since $s_A(2(K-1)) = s_B(2(K-1)) = K-1$, we obtain
%
\begin{eqnarray*}
   \Pr\{T_{2(K-1)}|(p_A,p_B,K)\} & \equiv & \theta_{ST}(K-1,K-1;p_A,p_B,K) \\  & = & \Pr\{B(K-1,p_A) + B(K-1,q_B) = K-1\}  \\
    & = & \left[\sum_{l=0}^{K-1}
{{K-1} \choose l} p_A^l q_A^{{K-1}-l} {{K-1} \choose {l}} q_B^{K-1-l} p_B^{l}\right].
\end{eqnarray*}

Finally, we note that if the ST reaches an STT, then we need to determine who serves first in this STT, which is indicated by $I_A(2K-1)$. Thus, we have 
%
\begin{eqnarray*}
\Pr\{\mbox{$A$ wins STT}|(p_A,p_B,K)\} & = & [\theta_{STT}(p_A,p_B)]^{I_A(2K-1)} [1-\theta_{STT}(p_B,p_A)]^{1 - I_A(2K-1)} \\ 
& = & [\theta_{STT}(p_A,p_B)]^{I_A(2K-1)} [1-\theta_{STT}(q_B,q_A)]^{1 - I_A(2K-1)}.   
\end{eqnarray*}

To establish $\theta_{ST}(p_A,p_B,K) = \theta_{ST}(q_B,q_A,K)$, it suffices to show that
%
$$\sum_{h=0}^{K-2} \Pr\{F_h|(p_A,p_B,K)\} = 
\sum_{h=0}^{K-2} \Pr\{F_h|(q_B,q_A,K)\}.$$
%
We establish this using an alternative way of looking at how a player could win based on looking at the complete paths of length $2(K-1)$ points.

Consider all paths or would-be paths of length $2(K-1)$. Note that some of these paths will never reach $2(K-1)$ points since it is possible that one player had already won $K$ points. However, think of extending the path until $2(K-1)$ points. Over these $2(K-1)$ points, each player will serve a total of $(K-1)$ times. 
Player $A$ then wins if
%
$\sum_{i=1}^{K-1} V_i + \sum_{j=1}^{K-1} (1-W_j) \ge K$
%
while player $B$ wins if
%
$\sum_{i=1}^{K-1} (1 - V_i) + \sum_{j=1}^{K-1} W_j \ge K.$
%
The other paths will lead to a tied score with each player having $(K-1)$ points. Under $(p_A,p_B)$ the probability of $A$ winning, without going to an STT, is
%
$$\Pr\{B(K-1,p_A) + B(K-1,q_B) \ge K\},$$
which is the alternative expression provided in the statement of the proposition for $\sum_{h=0}^{K-2} \Pr\{F_h|(p_A,p_B,K)\}$.
%
But note that this is also the probability of $A$ winning, without going to an STT, 
{\em under} $(q_B,q_A)$. As such, we have established, via an alternative viewpoint, that 
%
$$\sum_{h=0}^{K-2} \Pr\{F_h|(p_A,p_B,K)\} = 
\sum_{h=0}^{K-2} \Pr\{F_h|(q_B,q_A,K)\}.$$
% 
We note that showing {\em directly} the equality above using the formula for $$\sum_{h=0}^{K-2} \Pr\{F_h|(p_A,p_B,K)\}$$ seems not a straightforward matter.

%}
   
    \OLDPROOF{ The first expression follows from the addition rule of probability and the decomposition of the event that $A$ wins the set tie-breaker.

    Let us now consider the event that $A$ wins with a $(K,h)$ score, where $h \le (K-2)$. In such a case a total of $K+h$ points were played, and it must be the case that $A$ wins on the $(K+h)$th point. $A$ served this point if $I_A(K+h) = 1$. Thus, the probability that $A$ won on the $(K+h)$the point is 
    %
    $$p_A^{I_A(K+h)} [1 - p_B]^{1 - I_A(k+h)}.$$
    %
    On the first $(K+h-1)$ points, the total number of times that $A$ served is $s_A(K+h-1)$, while the number of $B$ serves is $s_B(K+h-1)$. The possible ways of getting a score of $(K,h)$ is for $A$ to win $j$ points on his/her serves, lose $s_A(K+h-1) - j$ of his/her serves, win $(K-1-j)$ of $B$'s serves, and lose $s_B(K+h-1) - ((K-1)-j)$ of $B$'s serves. This could happen for $j=0,1,2,\ldots,K-1$, with the understanding that the event has zero probability if $j > s_A(K+h-1)$ or $(K-1-j) > s_B(K+h-1)$. The number of possible configurations in which $A$ wins $j$ points on his/her serves is
    %
    $${s_A(K+h-1) \choose j} {s_B(K+h-1) \choose {(K-1) - j}}$$
    %
    and each of these configurations has probability
    %
    $$\left[p_A^j q_A^{s_A(K+h-1) - j}\right]
        \left[q_B^{(K-1)-j} p_A^{s_B(K+h-1) - ((K-1)-j)}\right].$$
        %
Summing all the probabilities for all the configurations results in the expression for $\Pr\{\mbox{$A$ wins on a $(K,h)$ score}\}$ in the statement of the proposition.

Finally, we have the last situation in which $A$ wins the tie-breaker, which is when both players reach $(K-1)$ points, which has probability (by the same analysis as above)
%
\begin{eqnarray*}
   \lefteqn{ \sum_{j=0}^{K-1} \left\{ {s_A(2K-2) \choose j} {s_B(2K-2) \choose {(K-1) - j}} \times \right. } \\ && \left.
        \left[p_A^j q_A^{s_A(2K-2) - j}\right]
        \left[q_B^{(K-1)-j} p_A^{s_B(2K-2) - ((K-1)-j)}\right]
        \right\}
\end{eqnarray*}
%
which then gets multiplied by the probability that $A$ wins the tie-breaker's tie-breaker. On the $[2(K-1)+1] = (2K -1)$th point, the server is indicated by $I_A(2K-1)$, so the probability that player $A$ will win the tie-breaker's tie-breaker is
%
$$[\theta_{STT}(p_A,p_B)]^{I_A(2K-1)} [1 - \theta_{STT}(p_B,p_A)]^{1 - I_A(2K-1)]},$$
%
since if player $B$ is serving on the $(2K-1)$th point, $B$ must lose the tie-breaker's tie-breaker, hence the component probability of $[1 - \theta_{STT}(p_B,p_A)]^{1 - I_A(2K-1)]}$ in the preceding equation.
%
Multiplying the two preceding displayed equations yield $\Pr\{(K-1,K-1)\ \mbox{score; and $A$ wins}\}$.

\RED{To establish the final identity in the statement of the proposition, first recall that $\theta_{STT}(p_A,p_B) = \theta_{STT}(q_B,q_A)$. Then we complete the proof by showing that
%
$$\sum_{h=0}^{K-2} \theta_{ST}(K,h;p_A,p_B,K) = \sum_{h=0}^{K-2} \theta_{ST}(K,h;q_B,q_A,K);$$
$$\theta_{ST}(K-1,K-1;p_A,p_B,K) = \theta_{ST}(K-1,K-1;q_B,q_A,K).$$

\BLUE{
Observe that $\sum_{h=0}^{K-2} \theta_{ST}(K,h;p_A,p_B,K)$ is the probability that $A$ wins the set tie-breaker without going into the set tie-breaker's tie-breaker (STT). Let $V_1, V_2, \ldots$ denote the independent Bernoulli random variables with parameter $p_A$ representing the outcomes when $A$ is serving the point, with $V_i = I\{\mbox{$A$ wins point}\}$. Similarly, let $W_1, W_2, \ldots$ denote the independent Bernoulli random variables with parameter $p_B$ representing the outcomes when $B$ is serving the point, with $W_j = I\{\mbox{$B$ wins point}\}$. Then, since over $2(K-1)$ points, each player will have $(K-1)$ service points, the event that $A$ wins the tie-breaker before an STT is equal to the event
%
$$\left\{ \sum_{i=1}^{K-1} V_i + \sum_{j=1}^{K-1} (1-W_j) \ge K\right\}.$$
%
Observe that $(1-W_j), j=1,2,\ldots$ are independent Bernoullis with parameter $q_B = 1 - p_B$. Thus,
%
\begin{eqnarray*}
\sum_{h=0}^{K-2} \theta_{ST}(K,h;p_A,p_B,K) & = &
\Pr\{BIN(K-1,p_A) + BIN(K-1,q_B) \ge K\} \\
& = & \Pr\{BIN(K-1,q_B) + BIN(K-1,p_A) \ge K\} \\
& = &\sum_{h=0}^{K-2} \theta_{ST}(K,h;q_B,q_A,K).
\end{eqnarray*}
%
This proves the first identity.

The second identity above follows immediately since
\begin{eqnarray*}
\lefteqn{ \theta_{ST}(K-1,K-1;p_A,p_B,K)  = } \\
& = &\Pr\{BIN(K-1,p_A) + BIN(K-1,q_B) =  K-1\} \\
& = & \Pr\{BIN(K-1,q_B) + BIN(K-1,p_A) = K-1\} \\
& = & \theta_{ST}(K-1,K-1;q_B,q_A,K).
\end{eqnarray*}
}}


}
\end{proof}

\end{document}
